\subsection{Separate charged-equilibrium study}
\label{sec:charged-extension}

Electromagnetism is introduced by the standard charged-boson-star EMKG action
\cite{JetzerVanDerBij1989,Pugliese2013,LopezAlcubierre2023},
\begin{equation}
S_{\rm EMKG}=\int\dd^4x\sqrt{-g}\left[
\frac{R}{16\pi G}-\frac14F_{\mu\nu}F^{\mu\nu}
-\frac12\sum_m\left(D_\mu\Phi_m^*D^\mu\Phi_m+
\mu^2|\Phi_m|^2\right)\right],
\qquad D_\mu=\nabla_\mu-\ii qA_\mu .
\label{eq:solved-emkg-action}
\end{equation}
For $A=A_t(r)\dd t$, define the gauge-invariant local frequency
$\Omega=\omega-qA_t$ and the radial electric-flux variable
\begin{equation}
Q_E=-\frac{r^2A_t'}{\alpha\gamma}.
\end{equation}
The Maxwell equations and the electrostatic contribution to the Einstein
equations give
\begin{align}
Q_E'&=\frac{q\kappa_\ell}{4\pi}\frac{\gamma}{\alpha}
r^2\Omega\psi^2,
&A_t'&=-\frac{\alpha\gamma Q_E}{r^2},\label{eq:charged-maxwell-pair}\\
M'&=\frac{\kappa_\ell r^2}{2}\left[
\frac{\psi'^2}{\gamma^2}+\left(\frac{\Omega^2}{\alpha^2}+
1+\frac{\ell(\ell+1)}{r^2}\right)\psi^2\right]
+\frac{2\pi Q_E^2}{r^2},\label{eq:charged-mass}\\
\frac{\alpha'}\alpha&=\gamma^2\left\{\frac Mr^2+
\frac{\kappa_\ell r}{2}\left[
\frac{\psi'^2}{\gamma^2}+\left(\frac{\Omega^2}{\alpha^2}-
1-\frac{\ell(\ell+1)}{r^2}\right)\psi^2\right]
-\frac{2\pi Q_E^2}{r^3}\right\}.
\label{eq:charged-lapse}
\end{align}
The Klein--Gordon equation is Eq.~\eqref{eq:kg} with $\omega$ replaced by
$\Omega(r)$.  Regularity imposes $Q_E(0)=0$.  The gauge is fixed by
$A_t(\infty)=0$, and the finite outer boundary is matched to
\begin{equation}
A_t=\frac{Q_E}{r},\qquad
\alpha^2=1-\frac{2M_{\rm ADM}}r+\frac{4\pi Q_E^2}{r^2},\qquad
M_{\rm ADM}=M(r)+\frac{2\pi Q_E^2}{r}.
\label{eq:rn-tail-convention}
\end{equation}
This normalization gives the physical Gauss charge
$Q_{\rm electric}=4\pi Q_E(\infty)=qQ_N$.

We solve the six-field nonlinear boundary-value problem and continue the
$a_1^0=0.08$ star through eleven charges up to $q=3.52$.  The Newtonian
balance value in this convention is $q_{\rm crit}=\sqrt{4\pi}=3.54491$.
Selected results appear in Table~\ref{tab:charged-sequence}.

\begin{table}[H]
\centering
\caption{Charged $\ell=1$ EMKG continuation at fixed $a_1^0=0.08$.  The
electric fraction is the integrated Maxwell-field energy divided by the ADM
mass.}
\label{tab:charged-sequence}
\small
\begin{tabular}{@{}ccccccc@{}}
\toprule
$q$ & $q/q_{\rm crit}$ & $\omega$ & $M_{\rm ADM}$ & $Q_{\rm electric}$ &
$C_{99}$ & $E_{\rm EM}/M$\\
\midrule
0.00 & 0.000 & 0.8518974 & 1.171271 & 0.0000 & 0.10488 & 0.000\\
1.00 & 0.282 & 0.8573647 & 1.235666 & 1.2771 & 0.10977 & 0.008\\
2.00 & 0.564 & 0.8760066 & 1.488137 & 3.0549 & 0.12868 & 0.038\\
3.00 & 0.846 & 0.9204041 & 2.366461 & 7.1672 & 0.19266 & 0.123\\
3.40 & 0.959 & 0.9602038 & 3.682903 & 12.4558 & 0.28644 & 0.225\\
3.52 & 0.993 & 0.9813001 & 4.725859 & 16.4221 & 0.35836 & 0.296\\
\bottomrule
\end{tabular}
\end{table}

At $q=3.52$, $\Omega(r)$ remains positive and reaches a minimum 0.34668 at
the center.  The Gauss identity closes with relative error
$5.94\times10^{-6}$, and the outer RN lapse and Coulomb-potential residuals
are $4.4\times10^{-16}$ and zero at stored precision.  At $q=3$, repeating
the solve with $r_{\max}=80,100,120$ produces full spreads
$3.30\times10^{-11}$ in $\omega$, $4.68\times10^{-10}$ in ADM mass,
$2.01\times10^{-9}$ in electric charge, and $2.46\times10^{-4}$ in $R_{99}$.
The exact $q=0$ solve reproduces the neutral mass to $2.1\times10^{-8}$ and
the neutral frequency to $6.5\times10^{-9}$.  Thus the electromagnetic
extension here is a solved backreacting equilibrium problem rather than the
force-balance estimate used only for physical-unit orientation.

\begin{figure}[H]
\centering
\includegraphics[width=0.96\textwidth]{figure_16_charged_quantum.pdf}
\caption{Controlled extensions.  Top panels: charged EMKG continuation in
$q/q_{\rm crit}$, showing the ADM mass, electric charge, compactness, and
Maxwell energy fraction.  Bottom panels: coherent-state axial-QNM dispersion
and mean bias across the boson-mass scaling.  The charged effect is nonlinear
and macroscopic near critical coupling; the specified coherent-state quantum
effect remains extremely small throughout the displayed range.}
\label{fig:charged-quantum}
\end{figure}

\subsection{Coherent-state occupation-noise toy estimate}
\label{sec:quantum-extension}

We define a restricted noise model rather than attaching an unspecified correction
to the classical eigenvalue.  The three $\ell=1$ magnetic modes are placed in
a product of equal-mean-occupation coherent states.  The stress tensor is
normal ordered relative to the asymptotic Minkowski vacuum.  Under this
renormalization prescription its leading expectation value is exactly the
classical multiplet stress used above, while the total occupation has Poisson
variance $\operatorname{Var}N=N$, the defining number-statistics property of a
coherent state \cite{Glauber1963}. The result therefore applies specifically
to this asymptotic-vacuum, normal-ordered coherent-state prescription;
vacuum-polarization models based on a different quantum state or subtraction
prescription define a different physical problem.

In the adiabatic one-mode approximation, occupation is proportional to the
computed Noether charge, $N\propto Q(a_1^0)$.  The local sequence gives
$Q/Q'=1.84320$ at $a_1^0=0.08$; this replaces the incorrect
$a_1^0/2=0.04$ coefficient implied by $a_1^0\propto\sqrt N$.  To propagate
the resulting fluctuations through the actual
axial spectrum, we solve two additional complex QNMs:
\begin{align}
\sigma(0.075)&=0.04740320289-6.11222\times10^{-7}\ii,\\
\sigma(0.085)&=0.05134426662-6.83394\times10^{-7}\ii.
\end{align}
Their normalized matching determinants are $4.08\times10^{-10}$ and
$5.23\times10^{-12}$.  Together with Eq.~\eqref{eq:first-axial-qnm}, centered
finite differences give
\begin{align}
\frac{\dd\sigma}{\dd a_1^0}&=0.39410637-7.21715\times10^{-6}\ii,\\
\frac{\dd^2\sigma}{\dd(a_1^0)^2}&=-1.9202877-4.41864\times10^{-4}\ii.
\label{eq:qnm-amplitude-sensitivity}
\end{align}
For Poisson occupation noise, expansion through $O(1/N)$ yields
\begin{align}
\Delta\sigma_{\rm rms}
&=\left|\frac{\dd\sigma}{\dd a_1^0}\right|
\left|\frac{Q}{Q'}\right|\frac{1}{\sqrt N},\label{eq:coherent-qnm-rms}\\
\langle\delta\sigma\rangle
&=\frac{1}{2N}\left[\sigma''\left(\frac{Q}{Q'}\right)^2
-\sigma'\frac{Q^2Q''}{(Q')^3}\right].
\label{eq:coherent-qnm-bias}
\end{align}
Across $10^{-22}\leq\mu/{\rm eV}\leq10^9$, the largest fractional RMS
dispersion is $1.09\times10^{-18}$ and the largest fractional mean bias is
$4.39\times10^{-36}$, both attained at the heaviest mass because it has the
smallest occupation.  At the stellar scaling $\mu=10^{-10}\,$eV they are
smaller by many additional orders.  This toy estimate gives a sharp answer
only for the stated occupation fluctuation: it does not measurably blur the
QNM in the displayed mass range.

This is not a calculation of $\langle T_{\mu\nu}\rangle_{\rm ren}$, loop
corrections, particle production, or quantum backreaction.  The conclusion is
deliberately model-specific.  A different quantum state,
a non-normal-ordered vacuum polarization, extra field species, or a chosen
higher-curvature effective action defines a different semiclassical theory
and is not numerically represented by Eqs.~\eqref{eq:coherent-qnm-rms}--
\eqref{eq:coherent-qnm-bias}.
